\paragraph{A.2. Оценка сумм}
\subparagraph{A.2.1}
Покажите, что сумма $\sum_{k=1}^{n} {1/k^2}$ ограничена сверху константой.
\begin{gather*}
\sum_{k=1}^{n} {1/k^2} \leq c \\
\sum_{k=1}^{n} {1/k^2} = 1 + \sum_{k=2}^{n} {1/k^2} \\
\sum_{k=2}^{n} {1/k^2} \leq \int_{1}^{n} {1/x^2} = 1 - \frac{1}{n} \\
\sum_{k=1}^{n} {1/k^2} \leq 1 + 1 - \frac{1}{n} = 2 - \frac{1}{n} \\ 
\sum_{k=1}^{n} {1/k^2} \leq 2
\end{gather*}

\subparagraph{A.2.2}
Покажите асимптотическую верхнюю границу суммы.
$$ \sum_{k=0}^{\lfloor \lg{n} \rfloor } {\lceil n/2^k \rceil}$$
\begin{gather*}
\sum_{k=0}^{\lfloor \lg{n} \rfloor } {\lceil n/2^k \rceil} \leq \sum_{k=0}^{\lfloor \lg{n} \rfloor }{n/2^k + 1} = \sum_{k=0}^{\lfloor \lg{n} \rfloor }{n/2^k} + \sum_{k=0}^{\lfloor \lg{n} \rfloor }{1} \\
a_1/a_0 = \frac{n * 2^k}{2^{k+1} * n} = \frac{1}{2} = r \\
\sum_{k=0}^{\lfloor \lg{n} \rfloor }{n/2^k} \leq a_0 * \frac{1}{1-r} = 2n
\sum_{k=0}^{\lfloor \lg{n} \rfloor } {\lceil n/2^k \rceil} \leq 2n + \lg{n} = O(n)
\end{gather*}

\subparagraph{A.2.3}
Применив разбиение ряда, покажите, что $n$-е гармоническое число представляет собой $\Omega(\lg{n})$.
\begin{gather*}
\sum_{k = 1}^{n} {\frac{1}{k}} \geq \sum_{i = 0}^{\left\lfloor \lg{n} \right\rfloor - 1} {\sum_{j = 0}^{2^i - 1} {\frac{1}{2^i + j}}} \geq \sum_{i = 0}^{\left\lfloor \lg{n} \right\rfloor - 1} {\sum_{j = 0}^{2^i - 1} {\frac{1}{2^{i+1}}}} = \\
= \sum_{i = 0}^{\left\lfloor \lg{n} \right\rfloor - 1} {\frac{1}{2}} \geq \frac{\lg{n} - 2}{2}  = \Omega(\lg{n})
\end{gather*}

\subparagraph{A.2.4}
Найдите приближенное значение $\sum_{k=1}^{n}{k^3}$ с помощью интегралов.

$\sum_{k=1}^{n}{k^3}$ - функция монотонно возрастающая
\begin{gather*}
\int_{0}^{n} {x^3dx} \leq \sum_{k=1}^{n}{k^3} \leq \int_{1}^{n+1} {x^3dx} \\
\int_{0}^{n} {x^3dx} = \frac{n^4}{4} \\
\int_{1}^{n+1} {x^3dx} = \frac{(n+1)^4}{4} - \frac{1}{4} \\
\frac{n^4}{4} \leq \sum_{k=1}^{n}{k^3} \leq \frac{(n+1)^4}{4} - \frac{1}{4} \\
\sum_{k=1}^{n}{k^3} = \Theta(n^4)
\end{gather*}

\subparagraph{A.2.5}
Почему для получения $n$-го гармонического числа мы не можем применить интегральное приближение непостредственно к сумме $\sum_{k=1}^{n} {\frac{1}{k}}$?

$\sum_{k=1}^{n} {\frac{1}{k}}$ - монотонно убывающая функция, тогда:
\begin{gather*}
\int_{1}^{n+1} {\frac{1}{x}dx} \leq \sum_{k=1}^{n} {\frac{1}{k}} \leq \int_{0}^{n} {\frac{1}{x}dx} \\
\end{gather*}
Интеграл $\int_{0}^{n} {\frac{1}{x}dx}$ является несобственным, равен бесконечности.
Следовательно, мы не можем оценить исследуемую функцию, используя интегральное приближение напрямую.